% !PRINTED 287
\weberpage{287}
%<<CONT
in agreement with the rule for the composition of two transformations in § 28. The substitution (11) is accordingly to be denoted by
\[
\begin{pmatrix} \partial, & c \\ b, & a \end{pmatrix}.
\]
The group formed by all substitutions of this form [according to (14)] we denote by $\mathfrak{L}$ (group of linear substitutions).

The functions (12) of $z'$, and therefore also the substitutions (13), remain unchanged if the four numbers $a$, $b$, $c$, $\partial$ are multiplied by one and the same factor not divisible by $p$, and from this there follows, by (10), the number of these functions, or the degree of the group $\mathfrak{L}$
\weq{15}{p(p^2-1),}
which one also finds easily by direct enumeration.

If in one of the substitutions (13) the four numbers $a$, $b$, $c$, $\partial$ are multiplied by a common factor indivisible by $p$, then the determinant $\Delta$ is multiplied by the \gesperrt{square} of this factor. Hence it is not the determinant $\Delta$, but rather its quadratic character, that is, the value of the symbol
\[
\left(\frac{\Delta}{p}\right)
\]
that is preserved under this multiplication.

If we compose two of the substitutions (13), their determinants are multiplied together, and from this it follows that \gesperrt{all substitutions of the group $\mathfrak{L}$ in which $\Delta$ is a quadratic residue of $p$ form a group among themselves, which we shall denote by $\mathfrak{L}_0$.}

The group $\mathfrak{L}_0$ is a (proper) divisor of the group $\mathfrak{L}$ of index 2.

If one composes the substitutions of $\mathfrak{L}_0$ with any substitution $(z, \beta z) = \begin{pmatrix} 1, & 0 \\ 0, & \beta^{-1} \end{pmatrix}$, where $\beta$, and therefore also $\beta^{-1}$, is a quadratic non-residue of $p$, then one obtains the whole group $\mathfrak{L}$.

If $\Delta$ is a quadratic residue of $p$, then one can choose a common factor by which to multiply $a$, $b$, $c$, $\partial$ in such a way that $\Delta \equiv 1 \ (\mathrm{mod}\ p)$, so that we can also represent the group $\mathfrak{L}_0$ by
%CONT>>
